\section{The persistent-state lifecycle}\label{sec:lifecycle}


\regime{Motion}{the operations that move state through the agent, and the five invariants that should govern them.}
The previous parts established what persistent state \emph{is} and what forms it takes. The next question is how it \emph{moves}: by what sequence of operations an observation becomes durable, governed, retrievable state, and by what later operations that state is revised, removed, accounted for, or undone. Memory, in the sense most of the literature uses the word, names only a few of these operations. An always-on agent observes, decides what to write, validates what it wrote, organizes it into something retrievable, retrieves it, acts on it, and then, once an outcome is known, updates it, forgets parts of it, audits how it was used, and on occasion rolls back the decisions a bad record caused. We call this sequence the persistent-state lifecycle. The unit of analysis is not a memory record alone but a typed transition over a set of records.

The lifecycle has two arcs, shown in Figure~\ref{fig:lifecycle}. The \emph{forward arc} runs observe, write, validate, organize, retrieve, act: it carries a signal from the environment into action. The \emph{return arc} runs update, forget, audit, rollback once the consequences of an action are known, reconciling state with reality. Plasticity, the capacity to absorb new information, lives mostly at write and update on the boundary between the arcs. Stability, the capacity to keep old commitments intact and reversible, is enforced at validate, audit, forget, and rollback. An agent that exposes the plastic operations without the stabilizing ones accumulates risk even as its task scores rise, a tension we develop in the part on continual adaptation and preview at the end of this one.

This framing differs from the standard memory pipeline in two places: the validate stage and the return arc. Record-centric memory taxonomies often leave both implicit, and the corpus quantifies that absence. Coding every work in the $N{=}435$ corpus by the lifecycle stages it exercises (a work may touch several) yields a sharply skewed profile: retrieve $269$, write $200$, act $141$, organize $128$, update $127$, audit $88$, validate $87$, observe $68$, forget $66$, and rollback $27$. The forward arc, especially retrieve and write, dominates. The return arc thins out, and rollback, the operation that closes the loop by repairing a state-affected decision, is the rarest capability in the corpus. The operations are therefore uneven in a structured way: the stages least represented in the coded corpus are the ones on which the invariants that distinguish a governed system from a memory-augmented one are established or checked. The stage-by-stage reading names representative mechanisms and benchmarks, contrasts them, and ends each stage on the governance gap it exposes.

\subsection{State records, transitions, and invariants}\label{subsec:records}

The six state axes (authority, scope, mutability, provenance, recoverability, actionability) become operational once we fix what a state item is and which changes to it are legal. We model a persistent state item as a typed record
\[
s = \langle v, a, c, m, p, r, k, \tau \rangle:
\]
a value $v$ (the remembered content); an authority $a$ naming what permitted $s$ to influence action (a user grant, a policy, a tool scope, a delegation); a scope $c$ over users, tools, time, and groups within which $s$ may be used; a mutability class $m$ (revisable, supersede-only, decaying, locked); a provenance chain $p$ recording the sources and transformations that produced $s$; a recoverability handle $r$ linking $s$ to the decisions and side effects it influenced; an actionability type $k$ drawn from a small closed set (evidence, preference, policy, skill, executable commitment); and a logical timestamp $\tau$. The agent's persistent state is the set of live records. The value $v$ is the only field most memory systems represent explicitly; the remaining seven are the governance envelope, and systems that discard the envelope cannot later recover the governance state it carried, however faithfully they preserve $v$.

Each lifecycle operation is then a typed transition over this set, and stating the transitions explicitly lets us name properties an always-on agent must actively preserve but an episodic agent satisfies vacuously by resetting. \textsc{Write} inserts a record with an authority and a provenance chain; in a governed system it inserts into a quarantined tier rather than the live set. \textsc{Validate} promotes a quarantined record to live, or rejects it, after checking that its authority is current and its provenance trusted. \textsc{Organize} re-represents live records (indexing, linking, summarizing, compressing) without changing which decisions they may license. \textsc{Retrieve} selects a subset of live records whose scope covers the current situation. \textsc{Act} conditions a decision on the retrieved subset and stamps the resulting action with the handles $r$ of the records that justified it. \textsc{Update} supersedes one record by a newer one within the same scope, preserving the old record's history rather than overwriting it. \textsc{Forget} tombstones a record and is obligated to propagate the tombstone to derived copies. \textsc{Audit} reconstructs, from $r$ and $p$, which records justified a given action. \textsc{Rollback} uses $r$ to revert the decisions and, where possible, the side effects that a now-removed or now-discredited record caused.

\subsubsection{The five invariants and the stages they govern}

Reading the transitions this way surfaces five invariants the loop should preserve. They are not a borrowed checklist; each is the property whose violation defines one of the persistence-specific failure classes treated later in the survey, and each is anchored to particular stages.

\emph{Authority monotonicity.} A record may influence an action only if a current, unrevoked authority covers it; authority may narrow over time but never silently re-broaden. An expired permission cannot license a write or a tool call. This invariant is enforced principally at validate (does this write carry a live authority?) and at act (is the authority that licensed this record still current at action time?), and it is exercised at update whenever a permission epoch advances.

\emph{Scope non-expansion.} A transition may narrow or preserve scope but never silently widen it. This is what blocks one user's preference from acting for another, or a fact learned in one task from being applied in an unrelated one. It governs organize (consolidation must not merge records across scopes into a wider-scoped summary) and retrieve (the retrieved set must not exceed the requester's scope).

\emph{Deletion propagation.} Tombstoning a record must also reach its summaries, embeddings, cached prompts, promoted tiers, and any derived records, so that erasure is a verifiable property of the store rather than an edit to a single retrieval index. It governs forget and is verified at audit.

\emph{Provenance preservation.} Consolidation may compress $v$ but must retain enough of the provenance chain $p$ to verify, attribute, and later revise the result. It governs organize and update, the two stages that transform $v$, and is the precondition for audit.

\emph{Rollback traceability.} Every action carries a handle back to the records that justified it, so that a later-discovered bad record yields a bounded, identifiable set of decisions and side effects to repair. It is established at act (the handle must be recorded) and consumed at rollback.

The mapping is summarized in Table~\ref{tab:invariant-stage}. The invariants split across the arcs unevenly: authority monotonicity and scope non-expansion are checked on the forward arc (at validate, retrieve, and act), while deletion propagation, rollback traceability, and the back half of provenance preservation all depend on the return arc, the same arc the corpus counts show to be sparse. Deletion propagation is the one invariant established wholly on the return arc (at forget, verified at audit); rollback traceability straddles the boundary, since the handle is set at act on the forward arc but consumed at rollback on the return arc, so a system that instruments act but never builds the return arc records the handle and then has nowhere to use it; and provenance preservation is at risk whenever its return-arc half, the update and audit steps, is the part a system omits. The invariants that most depend on the return arc are therefore the ones with the fewest stage opportunities in our coded corpus. Figure~\ref{fig:invariants} makes this proxy view concrete by counting, per taxonomy part, the works that touch each invariant's governing stage. It should not be read as direct evidence that those works test or enforce the invariant; it shows where the stages that would make such enforcement possible are present.

\begin{table}[t]
\centering
\small
\setlength{\tabcolsep}{4pt}\renewcommand{\arraystretch}{1.2}
\begin{tabularx}{\linewidth}{p{0.24\linewidth}Yp{0.20\linewidth}}
\toprule
Invariant & Governing stage(s) & Arc \\
\midrule
Authority monotonicity & validate, act (update) & forward \\
Scope non-expansion & organize, retrieve & forward \\
Provenance preservation & organize, update & both \\
Deletion propagation & forget (verified at audit) & return \\
Rollback traceability & act (set), rollback (use) & return \\
\bottomrule
\end{tabularx}
\caption{The five invariants and the lifecycle stages that establish or check them. Authority and scope are guarded on the forward arc; deletion propagation and rollback traceability live on the return arc, the arc the corpus shows is least implemented.}
\label{tab:invariant-stage}
\end{table}

\subsection{The forward arc: observe, write, validate, organize, retrieve, act}\label{subsec:forward}

\subsubsection{Observe and write: deciding what becomes state}

The forward arc begins before any memory is written, when the agent decides which observations are worth keeping. The corpus is thinnest here ($68$ works touch observe), largely because most memory systems treat the write decision as a fixed extraction policy rather than a governed choice. The dominant pattern is fact extraction and consolidation across sessions, exemplified by production-oriented systems that distill conversational turns into compact facts and merge them with prior knowledge \citep{chhikara2025mem0buildingproductionreadyai}. These systems are engineered for write throughput and retrieval recall, treating what to keep as a question of compression quality. The broader surveys of agent memory codify this view, organizing the field around extraction, distillation, structuring, and ranking, and making clear that write is studied as a content operation, not as the point where authority and provenance should first attach to a record \citep{zhang2024memorymechanism,du2026memoryautonomous,hu2025memoryageaiagents}. A recent survey reframing memory around a read/write/update lifecycle moves closer to the operational view we take, but still reasons at the level of mechanisms rather than typed transitions with governance fields \citep{huang2026rethinkingmemorymechanisms}.

A small number of works reveal what a richer write stage requires. The episodic-future-thinking memory of \citet{guo2026tmem} rehearses memories at write time in anticipation of future contexts rather than simply archiving them, so the write decision is conditioned on expected later use; this is a step toward write as a forward-looking governance choice, though it still lacks an authority field on the written record. A different and underappreciated move takes the write decision out of the agent's hands entirely: the memory-sandbox system turns an agent's memories into manipulable visual objects so a user can see, edit, delete, and share what is remembered, relocating write, organize, and forget from the agent to the human \citep{huang2023memorysandbox}. As an interface result, this matters, but it also exposes the gap: when write is a manual edit, there is no formal account of which edits are legal, so a user can introduce an inconsistent state with no validation gate to catch it.

The gap at observe and write is that the coded corpus offers no widely used criterion for what to encode versus discard that incorporates the governance envelope. Systems decide what to keep by salience or compression value, not by whether the record can carry an authority, a scope, and a provenance chain that later stages will need. A write that captures a useful value but no authority is a write that no downstream stage can govern.

\subsubsection{Validate: the missing gate}

Validate should sit between a write and the live state, promoting a record only after checking that its authority is current and its provenance trusted. The corpus shows it is barely a stage at all ($87$ works, most addressing it implicitly through retrieval filtering rather than as a distinct promotion gate). The staleness literature states the need directly: \citet{chao2026stalellmagentsknow} build a benchmark testing whether agents detect when stored memories have become stale or invalid, and find that agents frequently act on obsolete state. This is a validate failure surfaced at retrieve time: the record was never marked as requiring revalidation, so nothing checks its currency before use. The benchmark is a single-stage probe, however, and does not connect staleness detection to the write-organize-retrieve chain that produced the stale record in the first place.

A direct critique of the field's missing abstractions appears in the position survey of \citet{zhou2026agentnative}, which asks whether current memory systems are agent-native or retrofitted from retrieval and database stacks. The diagnosis is that without native abstractions, validate has nowhere to live: a vector store promotes whatever is written, because promotion is not a concept it represents. The survey identifies the gap but stops short of proposing the validate-stage invariants (authority monotonicity at promotion time) that would make a system demonstrably agent-native. The validate gap is central on the forward arc because this is where authority monotonicity should be enforced. Almost no system enforces it, so a record written under a since-revoked permission, or distilled from an untrusted tool output, can flow into live state and later license an action.

\subsubsection{Organize: consolidation that preserves or destroys governance}

Once written and (ideally) validated, records are organized into retrievable structure: indexed, linked, summarized, and compressed. This stage is comparatively well served ($128$ works) and contains many of the field's most developed mechanisms, but it is also where provenance preservation and scope non-expansion are most often quietly violated. The mechanisms divide into two families. The first organizes by \emph{summarization}: recursive hierarchical summarization condenses long dialogue history into semantic summaries while maintaining retrieval accuracy across sessions \citep{wang2023recursivesummary}. The second organizes by \emph{structure}: agentic memory with self-organizing, dynamically linked notes that evolve and refactor as new information arrives \citep{xu2025amemagenticmemoryllm}, and temporal knowledge graphs that organize conversational memory as timestamped facts with explicit validity windows for multi-hop temporal queries \citep{rasmussen2025zep}.

Contrasting these families exposes the governance cost of consolidation. Summarization buys compactness at the price of provenance: \citet{wang2023recursivesummary} does not formalize what information loss is acceptable, and a summary that scores well on aggregate recall can silently drop the source attributions audit and rollback later need. Structural organization buys multi-hop reasoning but introduces a scope hazard: the self-organizing links of \citet{xu2025amemagenticmemoryllm} can connect records across scopes, so a traversal can surface a fact outside the requester's authority, an inference channel the system does not guard. Temporal graphs are comparatively better on provenance because validity windows preserve when a fact held \citep{rasmussen2025zep}, yet even there the schema is not validated to preserve full lifecycle traceability under mutation. The gap at organize is that consolidation is optimized for retrieval quality and storage cost, not for the invariants it must not break: a consolidation step that compresses provenance away or merges scopes is a governance loss no later stage can recover.

\subsubsection{Retrieve: the most-studied stage, and what it still misses}

Retrieve is the field's center of gravity ($269$ works, more than any other stage), and the benchmarks here are mature. Long-term conversational memory benchmarks such as LoCoMo test whether agents retrieve semantically distant yet critical memories for complex multi-turn reasoning \citep{maharana2024evaluatinglongtermconversationalmemory}, and broader assistant benchmarks probe long-term interactive recall across many sessions \citep{wu2025longmemevalbenchmarkingchatassistants}. More recent work makes retrieval evaluation harder and more diagnostic. MemoryArena evaluates memory across interdependent multi-session tasks, where earlier actions and feedback should change later retrievals, with separate metrics for recall, induction, and contextual grounding \citep{he2026memoryarena}; incremental multi-turn evaluations decompose retrieval competence into distinct skills \citep{hu2026evaluatingmemoryllmagents}; and benchmarks of fine-grained relational discrimination test whether agents can tell complementary, nuanced, and contradictory memories apart, separating preservation, retrieval, and reasoning failures \citep{wang2026subtlememory}. Realistic, temporally evolving dialogue generators push retrieval evaluation toward heterogeneity and drift \citep{zhang2026rhelm}.

The limitation of this family, visible across these benchmarks, is that retrieval is usually scored against the value field $v$ alone. They measure whether the right content comes back, not whether the retrieved set respects scope, whether the records are still authorized, or whether a retrieved fact will correctly ground the downstream action. MemoryArena comes closest to the action-grounding question because its tasks are interdependent, but as its own analysis notes, it measures retrieval accuracy without diagnosing why retrieved memory fails to ground action or whether organization enabled the retrieval under interference \citep{he2026memoryarena}. The gap at retrieve is thus not retrieval quality, which is well advanced, but the absence of scope non-expansion as a scored property: a retrieval that surfaces a correct value from outside the requester's scope passes every benchmark in this family while violating the invariant.

\subsubsection{Act: binding state to consequence}

Act is where retrieved state conditions a decision with real consequences, and where the forward arc must stamp the action with the recoverability handles the return arc will need. The foundational mechanism is the interleaving of reasoning and action that lets an agent bind internal state to tool calls and validate intermediate outcomes \citep{yao2023reactsynergizingreasoningacting}; this established how agents act on state but said nothing about how persistent memory should guide action selection or how a failed action should update the state that licensed it. The benchmarks that put state-grounded action under stress are the stateful, tool-mediated environments. AppWorld tests whether agents track world-state changes across interactions over stateful apps and simulated people \citep{trivedi-etal-2024-appworld}; $\tau$-bench evaluates tool-agent-user interaction with explicit policy-compliance checks, so an action that violates a stated policy fails even when it completes \citep{yao2024taubenchbenchmarktoolagentuserinteraction}; and GAIA stresses multi-step tool use on real-world tasks \citep{mialon2023gaia}. The most recent, Momento, is pointed for our argument: it requires consequential tool-mediated actions over evolving goals across sessions, so an agent must treat prior-session history as state that may need re-validation before it acts, not as settled fact \citep{merin2026momento}.

These benchmarks are the closest the field comes to scoring the act stage as a state operation rather than an answer, and $\tau$-bench's policy compliance and Momento's re-validation requirement are genuine steps toward authority monotonicity at action time. The remaining gap is twofold. First, as the AppWorld analysis notes, when an action fails it is hard to attribute the failure to an incomplete write, a poor organization, or a weak retrieval, because the trajectory does not carry the per-record handles that would localize it \citep{trivedi-etal-2024-appworld}. Second, none of these benchmarks scores whether the action recorded a rollback handle: an action with an irreversible side effect (a payment, a deletion, an external message) is treated as correct if it satisfies the task, even if the state that licensed it later proves stale or unauthorized and the side effect cannot be traced back for repair. Act, in current practice, conditions on state but does not instrument the consequence, which is what makes the return arc hard to implement.

\subsection{The return arc: update, forget, audit, rollback}\label{subsec:return}

The return arc runs once an outcome is known, reconciling state with reality. The corpus counts collapse here: update $127$, audit $88$, forget $66$, rollback $27$. Update is comparatively healthy because revising a fact resembles a write. The governance-bearing operations, those that remove state, account for its use, and undo its effects, are progressively sparser, and rollback is near-absent. The pattern is simple: the field can change state more readily than it can remove, account for, or reverse it.

\subsubsection{Update: revision that is studied, propagation that is not}

Update supersedes a record with a newer one in the same scope. It is the best-served return-arc stage because revising a stored fact is mechanically close to writing one, and the conversational-memory and temporal-graph systems discussed under organize already implement it through validity windows and self-revising notes \citep{rasmussen2025zep,xu2025amemagenticmemoryllm}. The hard part of update is not changing the focal record but propagating the change to every record derived from it, and detecting when an update should have fired but did not. The staleness benchmark of \citet{chao2026stalellmagentsknow} is, read on the return arc, an update-propagation diagnostic: the obsolete record was never superseded, so the update that should have fired never did. Belief-drift studies extend this to softer mutation, measuring whether an agent's stored opinions remain temporally consistent or drift under repeated interaction \citep{myakala2026beliefshiftbenchmarkingtemporalbelief}. The gap at update is that supersession is treated as a local edit: superseding a fact rarely re-examines the summaries, skills, or downstream commitments derived from the old value, so a corrected fact can coexist with stale derivations of itself, a provenance-preservation failure the next stages are supposed to catch and usually cannot.

\subsubsection{Forget: deletion as a retrieval edit, not an erasure}

Forget is where deletion propagation must hold, and it is one of the least-implemented stages ($66$ works). The dominant treatment is decay rather than deletion: Ebbinghaus-inspired forgetting downweights records over time so less-used memories fade \citep{zhong2023memorybankenhancinglargelanguage}. Decay is attractive because it is cheap and graceful, but it does not satisfy a deletion request: a downweighted record is still present and still retrievable under the right query, precisely the failure mode a user invoking a right-to-be-forgotten would care about. The interface alternative, the memory-sandbox, lets a user explicitly delete a memory object \citep{huang2023memorysandbox}, closer to true erasure, but it deletes the object the user can see and does not propagate to summaries, embeddings, or promoted tiers derived from it. The gap at forget is the deletion-propagation invariant itself: in both the decay and the manual-edit treatments, forgetting edits one representation of a record while its derived copies survive, so erasure becomes a property of a single index rather than of the store. A fact a user asked to delete can persist in a summary computed before the deletion, and nothing in these mechanisms detects or repairs that.

\subsubsection{Audit: accounting for how state was used}

Audit reconstructs which records justified a given action, and it is sparse ($88$ works) and shallow where it exists. The interface line again offers the most concrete instance: the memory-sandbox makes the agent's state observable and lets a user trace how a memory edit changes subsequent behavior, audit in an interactive, human-driven form \citep{huang2023memorysandbox}. This reveals that audit requires the provenance chain $p$ and the recoverability handles $r$ to have been recorded on the forward arc; when those fields are absent, audit can only inspect the current value of state, not the lineage of a past action. The agent-native critique sharpens the point: a retrofitted retrieval stack has no place to store the per-action justification trail, so audit is reduced to after-the-fact guessing about which retrieved record caused a behavior \citep{zhou2026agentnative}. The gap at audit is structural rather than algorithmic: without provenance preservation upstream, there is no lineage to audit, which is why audit and rollback co-fail.

\subsubsection{Rollback: the rarest capability}

Rollback closes the loop. Given a record discovered to be bad, stale, poisoned, or unauthorized, rollback uses the handles $r$ to revert the decisions and side effects it caused. It is the rarest stage in the entire corpus: $27$ of $435$ works expose any rollback mechanism, and within lifecycle-focused work none directly implements it. The benchmarks that most need it, the stateful action environments, do not score it: AppWorld and $\tau$-bench verify whether an action was correct, not whether a wrong action licensed by stale state could be traced and undone \citep{trivedi-etal-2024-appworld,yao2024taubenchbenchmarktoolagentuserinteraction}, and Momento's re-validation requirement is a forward-arc guard against acting on stale state, not a return-arc mechanism for repairing an action already taken \citep{merin2026momento}. The OS-inspired memory systems discussed next provide the closest substrate for rollback, through checkpointing and paging, yet as we note they do not connect the substrate to the state-affected decisions rollback must repair. One caveat on the aggregate count: it sums three mechanically distinct operations the literature treats separately. Internal-state rollback reverts a stored fact or a consolidation decision; workflow rollback compensates a partially executed multi-step task; and external-effect rollback undoes an already-committed side effect such as a payment or a deletion. The aggregate of twenty-seven is therefore conservative for the easiest subclass and almost certainly an overcount for the hardest. External-effect rollback, the subclass that matters most because irreversible effects cannot be repaired by any internal mechanism, is rarer still.

Rollback is the survey's central lifecycle gap. Rollback traceability is established at act (the handle must be recorded) and consumed at rollback, so a field that does not instrument act for consequence cannot implement rollback even in principle. The practical consequence is that an always-on agent that takes an irreversible action on the basis of a record it later discovers was poisoned or unauthorized has, in the vast majority of the systems we surveyed, no bounded set of decisions to repair and no path back to a known-good state. This operation distinguishes a governed persistent-state system from a memory-augmented one, and on the corpus evidence it is the one the literature has built least.

\subsection{Operating-system abstractions for the lifecycle}\label{subsec:os}

A distinct line of work treats the whole lifecycle as a resource-management problem and imports operating-system abstractions to solve it. MemGPT pioneered this view, treating agent memory as OS-style paged storage with a context manager that externalizes long-term state and manages a bounded working set, decoupling the task horizon from the context budget \citep{packer2024memgptllmsoperatingsystems}. MemoryOS extends the analogy with explicit kernel-level paging, scheduling, and eviction policies and a hierarchical working-set bound for long-horizon reasoning \citep{kang2025memoryosaiagent}. A complementary mechanism precomputes and consolidates context offline, before a query arrives, to amortize inference compute, in effect scheduling the organize stage during idle time \citep{lin2025sleeptimecompute}.

The OS framing is lifecycle-aware because it makes transitions between tiers (working set to long-term store and back) explicit operations rather than implicit side effects, and it provides the checkpoint and restore primitives rollback would build on. Its limitation is that it governs the lifecycle for \emph{resource} correctness, not \emph{state} correctness. MemGPT establishes no principled criterion for which records should cross the working-set boundary beyond recency and capacity, and does not validate that paging preserves state fidelity across episodes \citep{packer2024memgptllmsoperatingsystems}. MemoryOS allocates and evicts efficiently but does not connect its eviction policy to the authority and provenance axes, so a record can be evicted before it is acted upon (a state-dependent failure) or a low-authority record promoted because it is frequently accessed \citep{kang2025memoryosaiagent}. The sleep-time consolidation work amortizes compute but provides no link from the precomputed state back to its source observations and no rollback if the anticipated context was wrong \citep{lin2025sleeptimecompute}. The gap is that the OS abstraction supplies useful \emph{machinery} for governance (tiers, scheduling, checkpoints) while leaving the governance \emph{policy} (which transitions preserve the five invariants) unspecified, the separation between mechanism and governance this survey argues the field must close.

\subsection{Lifecycle fidelity and the return-arc gap}\label{subsec:fidelity}

A final cluster of work does not implement a stage but tests whether the chain as a whole preserves fidelity across multi-session reuse, and these works are the most direct evidence for the survey's thesis. Stale-state detection \citep{chao2026stalellmagentsknow}, the agent-native critique of retrofitted abstractions \citep{zhou2026agentnative}, and the human-in-the-loop manipulation of in-memory state \citep{huang2023memorysandbox} together show that lifecycle dependencies are real and largely ungoverned: a write decision affects downstream organization fidelity, an organization choice affects retrieval under interference, and a missing validate gate lets a stale or untrusted record flow all the way to action. None provides a formalization of the invariants the chain must preserve or an automated mechanism to recover when the chain breaks; they reveal the dependencies without governing them.

Table~\ref{tab:stage-coverage} consolidates the per-stage picture developed above. The shape of the table is the argument. Coverage is heavy on the forward arc and concentrated at retrieve and write; it thins monotonically across the return arc to rollback. The invariant most exposed at each stage is governed, if at all, by a stage that is thin in the coded corpus: scope non-expansion depends on a retrieve and organize discipline retrieval benchmarks do not score, deletion propagation depends on a forget operation decay-based systems do not perform, and rollback traceability depends on instrumenting act for consequence and on a rollback stage that $27$ of $435$ works implement.

\begin{table}[t]
\centering
\small
\setlength{\tabcolsep}{4pt}\renewcommand{\arraystretch}{1.2}
\begin{tabularx}{\linewidth}{p{0.13\linewidth}p{0.06\linewidth}YY}
\toprule
Stage & Works & Representative mechanism / benchmark & Per-stage gap \\
\midrule
Observe & 68 & Anticipatory write rehearsal \citep{guo2026tmem}; user-driven write \citep{huang2023memorysandbox} & No governed criterion for what to encode; write captures $v$ without authority \\
Write & 200 & Fact extraction and consolidation \citep{chhikara2025mem0buildingproductionreadyai}; memory surveys \citep{zhang2024memorymechanism,du2026memoryautonomous} & Write studied as compression, not as where authority and provenance attach \\
Validate & 87 & Staleness detection \citep{chao2026stalellmagentsknow}; agent-native critique \citep{zhou2026agentnative} & The missing promotion gate; authority monotonicity unenforced \\
Organize & 128 & Recursive summary \citep{wang2023recursivesummary}; self-organizing notes \citep{xu2025amemagenticmemoryllm}; temporal graph \citep{rasmussen2025zep} & Consolidation drops provenance or merges scopes \\
Retrieve & 269 & LoCoMo \citep{maharana2024evaluatinglongtermconversationalmemory}; MemoryArena \citep{he2026memoryarena}; relational discrimination \citep{wang2026subtlememory} & Scored on value only; scope non-expansion unscored \\
Act & 141 & ReAct \citep{yao2023reactsynergizingreasoningacting}; AppWorld \citep{trivedi-etal-2024-appworld}; $\tau$-bench \citep{yao2024taubenchbenchmarktoolagentuserinteraction}; Momento \citep{merin2026momento} & Conditions on state but records no rollback handle for side effects \\
Update & 127 & Temporal validity \citep{rasmussen2025zep}; belief drift \citep{myakala2026beliefshiftbenchmarkingtemporalbelief} & Supersession is local; derived records not re-examined \\
Forget & 66 & Decay-based forgetting \citep{zhong2023memorybankenhancinglargelanguage}; manual delete \citep{huang2023memorysandbox} & Deletion edits one index; does not propagate to derived copies \\
Audit & 88 & Observable state manipulation \citep{huang2023memorysandbox} & No upstream provenance, so no lineage to audit \\
Rollback & 27 & OS-style checkpoint substrate \citep{packer2024memgptllmsoperatingsystems,kang2025memoryosaiagent} & No link from substrate to state-affected decisions; rarest stage \\
\bottomrule
\end{tabularx}
\caption{Stage coverage across the $N{=}435$ corpus, with a representative mechanism or benchmark and the governance gap each stage exposes. Coverage is heavy on the forward arc (retrieve, write) and thins monotonically across the return arc to rollback. The per-stage gap column shows that each return-arc invariant is governed, if at all, by a stage that is thin in the coded corpus.}
\label{tab:stage-coverage}
\end{table}

The persistent-state lifecycle is not a memory pipeline with two extra boxes. It is a loop whose forward arc is much more developed in the works we coded than its return arc, and the invariants that distinguish a governed persistent-state system from a memory-augmented agent are the ones the return arc enforces. Several open problems follow directly. The corpus contains few validate stages that enforce authority monotonicity at promotion time, so untrusted or unauthorized writes can flow into live state. It contains few consolidation disciplines at organize that preserve provenance and refuse scope-widening merges. It contains few forget operations that propagate deletion to derived state, few audit trails rooted in upstream provenance, and, most acutely, few rollback mechanisms that instrument action for consequence so a bad record yields a bounded set of decisions to repair. In the works surveyed here, the five invariants are more often implicit desiderata than benchmarked properties of the write-organize-retrieve-act chain. These gaps are the operational content of the claim that always-on agents are persistent-state systems whose state must be governed, and they lead directly to continual adaptation, where accumulated state compounds risk rather than competence when the return arc is missing.
